\section{Limitations}
\label{sec:limitations}

The first limitation is the instrument. Labels come from comparing a call with a benchmark's ground truth, which is a proxy for whether the action would have been harmful, and the audit of \cref{app:audit} shows that a parser can manufacture a family effect. A hand-checked sample of twenty-five items per run put the label error between zero and one in five on the JSON dialects before the repair and near zero on the XML dialects, and \cref{app:audit} reports the error measured on the repaired labels; what residual error remains inflates probes, because a mislabelled correct call has a surface form a probe can learn and confidence cannot, so it would move the result toward the side where internals appear needed.

The second is scale and breadth. Every checkpoint is at most three billion parameters and every family is represented by one or two checkpoints, so the two sides hold three checkpoints each, which is below what any test can separate, and the property that divides them is a hypothesis. Agents in production run on models of eight billion parameters and above, and \citet{yeats2026calls} find probe accuracy rising with size; whether confidence rises with it is the open question the matched pairs of \cref{sec:discussion} would answer.

The third is realism. Calls are single-step on three benchmarks of English requests, and the multi-turn evaluation replays the benchmark's own earlier calls rather than the model's trajectory with executed tool results, so it shows that a judge is not disturbed by a wrong call in the history and cannot show what a wrong tool result does. Cost is measured on one consumer GPU and not on a serving stack.

The last is direction. The two protocol choices of \cref{sec:controls} that move the numbers, the training population and the sign-fixing population, both moved them toward confidence, and the confidence summary chosen on training folds moves the Llama advantage down by \bestConfShrinkMax{} at most; nothing we found moves the Qwen3, Qwen3.5 or MiniCPM5 advantage up. The result should therefore be read as a lower bound on how often confidence suffices and an upper bound on where internals are needed.
